\SourceLeaf{325}%
und in den Ergebnissen nicht mehr vorkommen, können hier jedoch ohne Weiteres für einen anderen Zweck verwendet werden.

Durch die beiden Punkte auf der Oberfläche des Ellipsoids, auf die sich die Aufgabe bezieht, ziehe man eine geodätische Linie von zunächst unbestimmter Ausdehnung und wähle darauf einen beliebigen Anfangspunkt. Das Stück dieser Linie vom Anfangspunkt bis zu einem unbestimmten Punkt werde mit \(u\) bezeichnet. Am letztgenannten Punkt wird die geodätische Linie in Richtung wachsender \(u\) und der Meridian von Norden nach Süden orientiert. Der Winkel, den die geodätische Linie dort mit dem Meridian bildet, heiße \(X\). Breite und Länge des unbestimmten Punkts seien \(Y\) und \(Z\). Ich nehme an, dass die Längen von Westen nach Osten und die Azimute \(X\) von Süden in Richtung Westen zunehmen. Bezeichnet man nun wie bisher die große Halbachse und die Exzentrizität der erzeugenden Ellipse mit \(a\) und \(e\), so gilt aus bekannten Gründen
\[
\begin{aligned}
\frac{\mathrm{d}Y}{\mathrm{d}u}&=-\frac{\cos X(1-ee\sin Y^2)^{\frac32}}{(1-ee)a}\quad .\ .\ .\ .\ .\ .\ .\ .\ .\ .\quad(1)\\
\frac{\mathrm{d}Z}{\mathrm{d}u}&=-\frac{\sin X(1-ee\sin Y^2)^{\frac12}}{a\cos Y}\quad .\ .\ .\ .\ .\ .\ .\ .\ .\ .\quad(2)
\end{aligned}
\]
Ferner ist nach einem bekannten Lehrsatz die Größe
\[
\frac{\sin X\cos Y}{\BareRoot{(1-ee\sin Y^2)}}
\]
für alle Punkte derselben geodätischen Linie konstant. Durch logarithmisches Differenzieren folgt daraus
\[
\operatorname{cotang}X\,\mathrm{d}X=(\operatorname{tang}Y-\frac{ee\cos Y\sin Y}{1-ee\sin Y^2})\mathrm{d}Y=\frac{(1-ee)\operatorname{tang}Y}{1-ee\sin Y^2}\mathord{.}\mathrm{d}Y
\]
und in Verbindung mit (1) daher
\[
\frac{\mathrm{d}X}{\mathrm{d}u}=-\frac{\sin X\operatorname{tang}Y(1-ee\sin Y^2)^{\frac12}}{a}\quad .\ .\ .\ .\ .\ .\ .\ .\ .\quad(3)
\]

\indent Wir wollen unsere Aufgabe jedoch allgemeiner fassen und
\[
\frac{\mathrm{d}X}{\mathrm{d}u}=x,\qquad\frac{\mathrm{d}Y}{\mathrm{d}u}=y,\qquad\frac{\mathrm{d}Z}{\mathrm{d}u}=z
\]
setzen, wobei wir zunächst nur voraussetzen, dass \(x,y,z\) beliebige gegebene Funktionen der beiden Variablen \(X\) und \(Y\) sind. Durch weiteres Differenzieren erhält man
